\clearpage
\fancyhead[L]{\small ROHDE. ÜBER DIE MASSEN DER PLANETEN.}
\SourceLeaf{488}
\begin{center}
{\large ROHDE. ÜBER DIE MASSEN DER PLANETEN.\par}
\medskip
\rule{0.55\linewidth}{0.4pt}\\[2pt]
von \textsc{Zach}. Monatliche Correspondenz. Band XII. S. 413 . . 419. Oktober 1805.\\[2pt]
\rule{0.55\linewidth}{0.4pt}
\end{center}
\begin{center}Gutachten eines anderen Rezensenten über die Abhandlung des Hauptmanns \textsc{Rohde} über die Massen der Planeten, aus einem Brief dieses Rezensenten.\end{center}

\noindent --- --- --- Die --- --- --- Verteidigung des Hauptmanns \textsc{Rohde}, die Sie mir freundlicherweise mitgeteilt haben, hat mir als Beispiel für die noch viel schlechtere Verteidigung einer schlechten Sache einigen Spaß gemacht. Allerdings bedaure ich doch, dass ich eine gewisse Mitschuld daran trage, einen so undankbaren Streit ausgelöst zu haben. Meine Hoffnung, dass eine Zurechtweisung Hauptmann \textsc{Rohde} wenigstens vorsichtiger machen würde, hat sich nicht erfüllt. Seine unüberlegte Antwort, die er so voreilig hat drucken lassen, ist wieder voll neuer Beweise seiner --- --- --- --- --- Ich möchte mich noch einmal mit Ihnen über diese Sache unterhalten, ohne Ihrer Antwort vorzugreifen (Sie haben sie wahrscheinlich schon abgeschlossen, und die Gründlichkeit und die tiefe Überzeugung, die überall aus jener Rezension sprechen, bürgen dafür, dass niemand sie besser vertreten kann als Sie selbst). Mein Hauptgrund ist, dass Sie wissen möchten, was ich von dieser Erwiderung auf Ihre Kritik halte.

Meine Ansicht über diesen Streit ist folgende. In Ihrer Rezension wollten Sie nicht bestreiten, dass \textsc{Rohde}s Methode zur Bestimmung der Planetenmassen in abstracto oder rein theoretisch betrachtet richtig und zuverlässig sein könnte. Sie haben behauptet und bewiesen,

1) dass \textsc{Rohde} sie bei Ceres und Pallas völlig falsch angewandt hat und dass seine Ergebnisse dafür schon deshalb überhaupt nichts wert sein können.

2) dass es beim heutigen Stand der praktischen Astronomie nicht möglich ist, die mittleren Abstände der Planeten so genau zu bestimmen, wie es für diese Methode nötig wäre, außer wenn man sie aus der Umlaufzeit unter Berücksichtigung ihrer Massen berechnet.

Eine Verteidigung wie die des Hauptmanns \textsc{Rohde} gegen 1) ist mir noch nie begegnet. Mit anderen Worten besagt sie Folgendes. \emph{„Da meine Methode, wie der Rezensent in 2) bewiesen hat, doch ganz unbrauchbar und ohne praktischen Wert ist, wäre es lächerlich gewesen, bei ihrer Anwendung noch besondere Sorgfalt aufzuwenden.“} Leider hat er dabei einen weiteren schülerhaften Fehler begangen. Er schließt, dass das Vorrücken der Tagundnachtgleichen in einem Jahr 0.01419 in Zeit beträgt, also in 1684 Tagen noch nicht 0.07 Tage. So etwas würde man einem Anfänger in der Astronomie nicht verzeihen. Das Vorrücken der Tagundnachtgleichen in einem Jahr beträgt 50$''$, und die Sonne braucht 0.014 Tage, um es einzuholen. Sie braucht also 0.07 Tage, um das Vorrücken von 1684 Tagen einzu\SourceLeaf{489}holen. Ceres braucht dafür aber mehr als 0.3 Tage, und um so viel ist der siderische Umlauf länger als der tropische.

Alles, was \textsc{Rohde} gegen 2) vorbringt, ist völlig nichtssagendes Geschwätz, das überhaupt nicht zur Sache gehört. Er meint, die Beobachtungskunst sei heute so weit fortgeschritten, dass die sechste Dezimalstelle der mittleren Abstände als zuverlässig gelten könne. Es ist jedoch allgemein bekannt, dass sämtliche Angaben über die mittleren Abstände der Planeten und des \textsc{Halley}schen Kometen in astronomischen Werken aus den Umlaufzeiten berechnet sind, entweder mit oder ohne Berücksichtigung der Massen. \textsc{Rohde} durfte also die Angaben in \textsc{La Place}s \emph{Exposition du Système du Monde} etc. gar nicht verwenden. Es ist daher reine Einbildung, wenn er behauptet, die Übereinstimmung seiner Ergebnisse für Planeten mit Trabanten mit den aus diesen Angaben gewonnenen Ergebnissen bestätige die Gültigkeit seiner Methode.

\textsc{Rohde} scheint übrigens selbst gemerkt zu haben, wie nichtig seine Antwort ist (das Gegenteil wäre auch unbegreiflich, sofern man nicht die seltsamsten Voraussetzungen annehmen wollte). Denn er versucht einzulenken und den Eindruck zu erwecken, er habe gar nicht vorgehabt, bessere Bestimmungen der Planetenmassen zu liefern. Er habe bloß eine streng richtige Methode gesucht, die jedem Einwand standhält, ohne sich darum zu kümmern, ob sich die nötigen Elemente und Daten heute schon mit hinreichender Genauigkeit aus Beobachtungen gewinnen lassen. Dazu sind nun drei Bemerkungen zu machen.

\noindent I) \emph{Diese Behauptung ist unwahr}, denn Herr \textsc{Rohde} behauptet am Ende seiner Abhandlung ganz deutlich, alle bisher für Merkur, Mars und Venus angestellten Störungsrechnungen seien wertlos. Zum Teil liege das daran, dass man ihre Massen ganz falsch angesetzt habe, und zum Teil daran, dass man die unvergleichlich großen Einwirkungen von Ceres, Pallas und Juno habe vernachlässigen müssen. Ferner habe er die Astronomen auf dem schnellen Weg des Drucks über seine wichtige Entdeckung informieren wollen und nicht auf dem langsamen Weg der Mitteilung im Manuskript. So sollten sich alle, die rechnen konnten, sogleich zusammentun, um die Störungsrechnungen anhand seiner neuen Massenangaben ganz umzugestalten, \emph{pour accélérer de tous côtés la récolte, à laquelle les forces d'un seul homme ne suffisent pas, et les miennes point du tout} (gewiss das Wahrste in der ganzen Abhandlung), \emph{vû la grande quantité de combinaisons analytiques nouvelles etc.}

\noindent II) Es gibt nur einen Weg, um den mittleren Abstand eines Planeten ohne Berechnung aus der Umlaufzeit zu bestimmen, oder um überhaupt irgendeinen Abstand aus Beobachtungen zu gewinnen. Man berechnet ihn trigonometrisch in dem allbekannten Dreieck, dessen drei Winkel die Kommutation, die Elongation und die jährliche Parallaxe sind und in dem man den Abstand der Sonne von der Erde als bekannte Seite annimmt. Die Winkel findet man aus der Länge der Sonne und der geozentrischen und heliozentrischen Länge des Planeten. Letztere muss man aus den Tafeln berechnen. Bekanntlich lässt sich aber für keinen Planeten die Genauigkeit dieser berechneten Länge auf 5$''$ und mehr verbürgen. Ebenso wenig hat man die beobachtete geozentrische Länge der Sonne und der Planeten bis auf ein paar Sekunden genau erhalten. Diese Werte sind sehr niedrig angesetzt. Daraus lassen sich für die verschiedenen Planeten leicht Ergebnisse ableiten, die zeigen, dass die dadurch entstehenden Fehler in den Abständen gegenüber der erforderlichen Genauigkeit enorm sind und \textsc{Rohde}s Methode ganz unbrauchbar machen.

Wird \textsc{Rohde} aus diesem Schlupfwinkel vertrieben, wird er sich damit begnügen zu sagen, seine Methode sei doch theoretisch richtig und lasse sich vielleicht nach Jahrtausenden anwenden, wenn die Beobachtungen zehn- oder hundertmal so genau seien wie heute. Dagegen bemerke ich nun

\SourceLeaf{490}
III) \textsc{Rohde}s Methode ist auch theoretisch betrachtet falsch, führt zu unvermeidlichen Zirkelschlüssen und ist deshalb ganz und gar unbrauchbar.

Ich habe bereits bemerkt, dass man den gleichzeitigen Abstand der Sonne von der Erde unbedingt braucht, um den Abstand eines Planeten aus Beobachtungen zu finden. Je genauer man diesen bestimmen will, desto genauer muss man jenen kennen. Daher müssen sowohl die Elemente der Sonnenbewegung als auch alle Störungen, die den Abstand der Sonne von der Erde beeinflussen, aufs Genaueste bekannt sein, also insbesondere auch die Störungen durch den Planeten selbst, dessen Abstand man bestimmen will. Also (und darin liegt die \emph{Pointe} der Sache) muss seine Masse schon bekannt sein. Ist sie es nicht, sind die Störungsgleichungen für den Abstand der Sonne von der Erde unvollständig, dann ist dieser Abstand falsch, dann ist der berechnete Abstand des Planeten von der Sonne falsch, dann ist auch der mittlere Abstand falsch, der auf einem oder mehreren wahren Abständen beruht. Ein Beispiel wird die Sache verdeutlichen.

Angenommen, jemand wollte nach \textsc{Rohde}s Methode die Masse des Mars bestimmen. Um seinen mittleren Abstand aus Beobachtungen zu bestimmen, bräuchte er einen oder mehrere wahre Abstände und dazu die Abstände der Sonne von der Erde. In unseren älteren Sonnentafeln fehlen aber alle von Mars abhängigen Gleichungen, weil seine Masse unbekannt ist (in Freiherrn von \textsc{Zach}s neuesten Sonnentafeln Arg. IV. und Arg. XV). Allein wegen dieses Mangels könnte der Abstand der Erde von der Sonne um $0.0000058 + 0.0000011 = 0.0000069$ fehlerhaft sein. Der daraus entstehende Fehler beim Abstand des Mars von der Sonne wäre in demselben Verhältnis größer, in dem dieser Abstand jenen übertrifft, und betrüge somit mehr als $0.00001$. Welchen Einfluss ein solcher Fehler im Abstand auf die Masse hätte, muss man Hauptmann \textsc{Rohde} selbst berechnen lassen. Auch schweige ich davon, dass man zur Bestimmung der Masse eines Planeten nach \textsc{Rohde}s Methode sowohl diese Masse schon kennen müsste als auch die aller anderen Planeten. Denn sonst könnte man weder eine genaue Theorie der Erde haben (um den Abstand der Sonne von der Erde zu berechnen) noch eine genaue Theorie des Planeten selbst. Seine heliozentrische Länge müsste man aus den Tafeln nehmen. Seine tatsächlichen Abstände könnte man zur Bestimmung des mittleren Abstands nur verwenden, wenn man sie vorher vom Einfluss aller fremden Störungen befreit und so rein elliptisch gemacht hätte. Doch genug und schon viel zu viel über dieses unreife Produkt.
